\documentclass[12pt]{article}
\usepackage[utf8]{inputenc}
\usepackage[T1]{fontenc}
\usepackage[margin=1in]{geometry}
\usepackage{setspace}
\usepackage{xcolor}
\usepackage[colorlinks=true, urlcolor=blue!60!black, linkcolor=blue!60!black]{hyperref}
\setlength{\parindent}{2em}
\setlength{\parskip}{1.25em}
\renewcommand{\baselinestretch}{1.0}
\begin{document}
\begin{center}
{\LARGE \textit{United States v. Vaello Madero}, 596 U.S. 159 (2022)\par}
\vspace{0.8em}
{\large Case Brief --- Four-Page Edition\par}
\end{center}
\vspace{1em}

\section{Memory Jogger}

The United States sued a disabled U.S. citizen to recover SSI payments made after he moved to Puerto Rico. The Court held 8--1 that the Fifth Amendment does not require Congress to extend SSI to Puerto Rico, because the territory's general exemption from federal income taxes supplies a rational basis under the Territory Clause (\href{https://www.law.cornell.edu/supremecourt/text/20-303}{596 U.S. 159}).

\section{Detailed Case Facts}

Jose Luis Vaello Madero, a U.S. citizen born in Puerto Rico in 1954, moved to New York in 1985. In 2012, after a severe illness, he began receiving Supplemental Security Income, a means-tested federal program enacted in 1972 for low-income aged, blind, and disabled people. By statute SSI is limited to residents of the 50 States and the District of Columbia (later extended by covenant to the Northern Mariana Islands); Puerto Rico residents receive instead a less generous, capped federal-territorial program, Aid to the Aged, Blind, and Disabled.

In 2013 Vaello Madero moved back to Puerto Rico. Unaware of the move, the Social Security Administration kept depositing monthly payments of \$733 to \$808 into his account.

\section{Procedural History}

In 2016, when he applied for Title~II benefits, the SSA learned of his move, terminated SSI retroactively, and in 2017 sued to recover \$28,081. He defended on the ground that excluding Puerto Rico residents violated the equal-protection component of the Fifth Amendment. The District of Puerto Rico granted him summary judgment, and the First Circuit unanimously affirmed, reasoning that SSI recipients are too poor to pay income taxes anywhere. The Supreme Court reversed on April 21, 2022.

\section{Judicial Votes}

Reversed, \textbf{8--1}.
\begin{itemize}
\item \textbf{Majority:} Justice Kavanaugh, joined by Chief Justice Roberts and Justices Thomas, Breyer, Alito, Kagan, Gorsuch, and Barrett.
\item \textbf{Concurring:} Justice Thomas; Justice Gorsuch.
\item \textbf{Dissenting:} Justice Sotomayor.
\end{itemize}

\section{Holding}

The Constitution does not require Congress, legislating under the Territory Clause, to extend SSI to residents of Puerto Rico. Under deferential rational-basis review, Puerto Rico's general exemption from most federal income, gift, estate, and excise taxes supplies a rational basis for distinguishing its residents in structuring benefits programs. Footnote~1 disclaimed any implication that Congress could exclude residents of individual States.

\section{Analysis of Opinions}

\textbf{Majority (Kavanaugh).} The Territory Clause gives Congress broad authority to make all needful rules for the territories, and since 1898 Congress has kept distinct tax and benefits regimes for Puerto Rico. Its residents pay Social Security, Medicare, and unemployment taxes and receive those benefits, but are generally exempt from federal income, gift, estate, and excise taxes on local income. \textit{Califano v.\ Torres} (upholding the SSI exclusion against a right-to-travel claim) and \textit{Harris v.\ Rosario} (upholding lower AFDC reimbursements) controlled. Congress may weigh the general balance of benefits and burdens without a dollar-to-dollar comparison. Requiring symmetry in benefits would invite demands for symmetry in taxes, burdening Puerto Rico's economy, and Congress remains free to extend SSI by statute.

\textbf{Thomas (concurring).} Thomas questioned the originalist basis of reverse incorporation under \textit{Bolling v.\ Sharpe} and \textit{Adarand}, suggesting that the Fourteenth Amendment's Citizenship Clause may better ground a ban on federal racial discrimination in civil rights.

\textbf{Gorsuch (concurring).} Because no party asked to overrule the \textit{Insular Cases}, Gorsuch joined, but he condemned them as a colonial invention grounded in racial stereotypes, mocking a regime in which uninhabited Palmyra Atoll is ``incorporated'' while three million citizens in Puerto Rico lack jury-trial rights, and urged their overruling.

\textbf{Sotomayor (dissenting).} The exclusion fails even rational-basis review. \textit{Torres} and \textit{Rosario} were summary rulings on different claims, resting on the false premise that Puerto Ricans pay no federal taxes, when they pay \$3--4 billion a year. SSI is a uniform, direct program for people too poor to owe income tax anywhere, so tying it to territorial tax status is ``antithetical to the entire premise of the program.'' The logic would let Congress exclude low-revenue States, and Puerto Ricans lack voting representation to fix the disparity. She agreed that the \textit{Insular Cases} should be overruled.

\section{Examples: Future Applications}

\begin{enumerate}
\item \textbf{Same-Side (Upheld): A Heating and Cooling Credit.} A new IRS-administered credit excludes residents of Guam and Puerto Rico, which keep separate block grants. Their distinct tax structures supply a rational basis.
\item \textbf{Same-Side (Upheld): Capped Territorial Medicaid.} Fixed block funding for territories, unlike the States' open-ended match, survives under the balance-of-burdens reasoning and \textit{Rosario}.
\item \textbf{Opposite-Side (Rationale Unavailable): Excluding a State.} Excluding West Virginians because they pay less income tax cannot rely on the Territory Clause, and Footnote~1 reserved the question.
\item \textbf{Opposite-Side (Strict Scrutiny): An Ancestry Line.} Excluding only residents of Hispanic ancestry is a racial or national-origin classification that fails strict scrutiny under \textit{Adarand}.
\item \textbf{Fence-Sitter: A Payroll-Funded Benefit.} A family-leave benefit funded by a payroll tax collected in Puerto Rico but paid only in the States removes the tax-exemption rationale, yet rational-basis review accepts any conceivable justification.
\end{enumerate}

\section{Critique}

\textbf{The fit problem.} The case shows how lenient rational-basis review operates at the level of aggregate jurisdiction rather than individual fit: no SSI-eligible New Yorker pays federal income tax either, yet the Court let Puerto Rico's territorial tax exemption stand in as a proxy for benefits eligibility.

\textbf{Footnote~1 and the \textit{Insular Cases}.} Footnote~1 is undertheorized: if benefits may track aggregate fiscal contributions, why not exclude low-revenue States, absent some unstated principle of equal state footing or political-process representation (cf.\ \textit{Carolene Products} Footnote~4)? And Gorsuch's concurrence exposes an irony: overruling the \textit{Insular Cases} would not have changed the result, because the majority assumed equal protection applied in full and still upheld the statute.

\section{Key Quotations}

\begin{quote}
``In devising tax and benefits programs, it is reasonable for Congress to take account of the general balance of benefits to and burdens on the residents of Puerto Rico. In doing so, Congress need not conduct a dollar-to-dollar comparison of how its tax and benefits programs apply in the States as compared to the Territories, either at the individual or collective level.''
\end{quote}
\noindent --- Justice Kavanaugh, Opinion of the Court, 596 U.S. at 165--66.

\begin{quote}
``It is past time to acknowledge the gravity of this error and admit what we know to be true: The \emph{Insular Cases} have no foundation in the Constitution and rest instead on racial stereotypes. They deserve no place in our law.''
\end{quote}
\noindent --- Justice Gorsuch, concurring, 596 U.S. at 180.

\begin{quote}
``In my view, there is no rational basis for Congress to treat needy citizens living anywhere in the United States so differently from others. To hold otherwise, as the Court does, is irrational and antithetical to the very nature of the SSI program and the equal protection of citizens guaranteed by the Constitution.''
\end{quote}
\noindent --- Justice Sotomayor, dissenting, 596 U.S. at 190.

\end{document}
